% ============================================================
%  Invoy VLM — Research Paper
%  Two-column, 10pt, standard article class (no special templates)
%  Compile: pdflatex (two passes) from the paper/ directory
%  Figures: ../eval_reports/three_model_plots/
% ============================================================
\documentclass[10pt, twocolumn]{article}

% --- geometry ---------------------------------------------------
\usepackage[top=1in, bottom=1in, left=0.75in, right=0.75in,
            columnsep=0.25in]{geometry}

% --- fonts & typography -----------------------------------------
\usepackage[T1]{fontenc}
\usepackage[utf8]{inputenc}
\usepackage{microtype}
\usepackage{lmodern}

% --- maths ------------------------------------------------------
\usepackage{amsmath}
\usepackage{amssymb}

% --- graphics ---------------------------------------------------
\usepackage{graphicx}
\graphicspath{{figures/}}

% --- tables -----------------------------------------------------
\usepackage{booktabs}
\usepackage{array}
\usepackage{multirow}

% --- captions ---------------------------------------------------
\usepackage{caption}
\captionsetup{font=small, labelfont=bf, skip=4pt}

% --- colour & links ---------------------------------------------
\usepackage{xcolor}
\definecolor{linkblue}{RGB}{0,70,140}
\usepackage[colorlinks=true,
            linkcolor=linkblue,
            citecolor=linkblue,
            urlcolor=linkblue,
            pdftitle={Invoy VLM: Benchmarking Local Vision-Language Models
                      for Continuous Desktop Activity Tracking},
            pdfauthor={Hamza Iqbal}]{hyperref}

% --- headers & footers ------------------------------------------
\usepackage{fancyhdr}
\pagestyle{fancy}
\fancyhf{}
\fancyhead[L]{\small\textit{Invoy VLM}}
\fancyhead[R]{\small\textit{Local VLMs for Desktop Activity Tracking}}
\fancyfoot[C]{\thepage}
\renewcommand{\headrulewidth}{0.4pt}
\setlength{\headheight}{13pt}

% --- section formatting -----------------------------------------
\usepackage{titlesec}
\titleformat{\section}{\normalfont\large\bfseries}{\thesection.}{0.5em}{}
\titleformat{\subsection}{\normalfont\normalsize\bfseries}{\thesubsection.}{0.5em}{}
\titlespacing*{\section}{0pt}{10pt plus 2pt minus 2pt}{4pt plus 1pt}
\titlespacing*{\subsection}{0pt}{6pt plus 2pt}{2pt plus 1pt}

% --- abstract formatting ----------------------------------------
\renewenvironment{abstract}{%
  \small\begin{center}{\bfseries Abstract}\end{center}%
  \begin{quotation}\noindent\ignorespaces
}{%
  \end{quotation}\medskip\hrule\medskip
}

% ================================================================
\title{%
  \vspace{-1.5em}%
  \LARGE\bfseries Invoy VLM: Benchmarking Local Vision-Language Models\\[4pt]
  for Continuous Desktop Activity Tracking%
  \vspace{0.3em}
}
\author{%
  Hamza Iqbal\\[2pt]
  \small \texttt{invoy\_sdk} Research Project\\
  \small April 2026
}
\date{}
% ================================================================

\begin{document}

\maketitle
\thispagestyle{fancy}

% ================================================================
\begin{abstract}
Continuous desktop activity tracking---automatically logging what a user is doing
from periodic screenshots---has practical applications in productivity monitoring,
workflow reconstruction, and personal analytics. We present \textbf{Invoy SDK}, a
lightweight system that captures screenshots every ten seconds and routes each frame
through a locally-resident Vision-Language Model~(VLM) to produce a natural-language
activity description and a change summary relative to the previous frame. We evaluate
three open-weight VLMs in this pipeline---Qwen2-VL~2B, Qwen2.5-VL~3B, and
Qwen2-VL~7B---over 136 frames drawn from a real Linux workstation session spanning
approximately 23~minutes. Using a suite of proxy-based evaluation metrics that require
no human annotation, we find: (1)~parameter count is the dominant predictor of
description richness; (2)~architectural improvements in the Qwen2.5-VL family do not
compensate for reduced parameter count at the 3B scale; and (3)~the 2B model achieves the highest Change-Delta F1 (0.334 vs.\ 0.200 for 7B),
meaning its change texts are the most lexically grounded in the tokens that actually
changed between activity frames---a reversal driven by the 7B model's tendency to
paraphrase rather than echo the activity delta.
\end{abstract}

% ================================================================
\section{Introduction}
\label{sec:intro}

Capturing \emph{what} a person is working on throughout a session has long relied on
application-level instrumentation---browser extensions, IDE plug-ins, OS-level hooks---
each covering only a narrow slice of activity and requiring explicit integration. A
purely visual approach, in which periodic desktop screenshots are interpreted by a
vision-language model, offers a unified, application-agnostic alternative.
Recent open-weight VLMs with strong visual comprehension have made local, on-device
inference practical at consumer GPU memory budgets, enabling privacy-preserving
deployments where no screenshot ever leaves the workstation.

Despite this promise, it is unclear how VLM quality---particularly at the compact
2--7 billion parameter range accessible on a single GPU---trades off against model
size for this specific task. Desktop screenshots are dense with small text, overlapping
windows, and domain-specific vocabulary; their interpretation requires both fine-grained
OCR-like reading and higher-level reasoning about what task a person is performing.

This paper addresses three research questions:
\begin{enumerate}
  \item[\textbf{RQ1.}] How does activity description quality scale with parameter count
        in the 2--7B range for this task?
  \item[\textbf{RQ2.}] Does the newer Qwen2.5-VL architecture outperform Qwen2-VL at a
        comparable or lower parameter count?
  \item[\textbf{RQ3.}] Can any of these models reliably detect task switches from
        screenshot pairs alone?
\end{enumerate}

We make the following contributions:
\begin{itemize}
  \item An open-source SDK (\textbf{Invoy}) for screenshot capture, local VLM inference,
        and structured SQLite logging, designed to make multi-model comparison
        reproducible.
  \item A proxy evaluation suite---including Change-Delta F1 and keyword-transition
        precision/recall---that quantifies activity description quality and change-text
        grounding without requiring human annotation.
  \item An empirical comparison of three Qwen-family VLMs on 136 real-world desktop
        frames, yielding actionable deployment guidance.
\end{itemize}

The remainder of the paper is structured as follows.
Section~\ref{sec:system} describes the Invoy system architecture.
Section~\ref{sec:eval} details the evaluation methodology.
Section~\ref{sec:results} presents results.
Section~\ref{sec:discussion} discusses findings and limitations.
Section~\ref{sec:related} surveys related work.
Section~\ref{sec:conclusion} concludes.

% ================================================================
\section{The Invoy System}
\label{sec:system}

Invoy consists of three loosely-coupled layers: a capture layer, an analysis layer, and
a storage layer. The design intentionally separates data collection from inference so
that the same screenshot archive can be re-analysed with different models or prompts
without re-running capture.

\subsection{Capture Layer}

A configurable capture loop (\texttt{collect\_periodic\_screenshots.py}) saves one
full-resolution PNG per interval (default: 10~s). The implementation uses the
\texttt{mss} library for X11 sessions and falls back to \texttt{gnome-screenshot} on
GNOME Wayland. Screenshots are named by UTC timestamp, providing a sortable archive
that is independent of any database.

\subsection{Analysis Layer}

\texttt{Qwen2VLBackend} (\texttt{invoy\_sdk/vlm\_backends.py}) wraps the Hugging Face
\texttt{transformers} library for Qwen2-VL~\cite{wang2024qwen2vl} and
Qwen2.5-VL~\cite{qwen25vl2025} model families. The correct model class is selected
automatically from the model-id string at load time.

Each frame triggers two separate inference calls:
\begin{enumerate}
  \item \textbf{Activity prompt.} The current screenshot alone is passed with a prompt
        asking the model to describe, in two to three sentences, what the person is
        currently doing on screen: which application is active, what file or page is
        open, and what task appears to be in progress.
  \item \textbf{Change prompt.} The activity text produced for the previous frame and the
        activity text just produced for the current frame are passed as plain text to the
        model. No image is used at this stage. The model is asked to describe what changed
        between the two descriptions and to state whether the task appears to be the same
        or different.
\end{enumerate}

\noindent
\textbf{Design rationale for text-only change detection.}
A natural alternative would be to pass both the previous and current screenshots
directly to the VLM as a two-image prompt, letting the model reason visually
about what changed. However, this was not considered feasible for the deployment
scenario Invoy targets. The central research question of this work is whether
these models can run \emph{locally and continuously} on a single consumer GPU at a
10-second capture interval. At that cadence, each activity inference call already
consumes several seconds of GPU time. Adding a second two-image inference call for
change detection would roughly double the per-frame compute cost, pushing the
effective loop period well beyond the 10-second target and making sustained
real-time operation impractical. Comparing the two activity \emph{texts}
instead---a pure language model call over short strings---adds negligible overhead
and keeps the pipeline within its resource budget. This choice is a direct
consequence of the feasibility constraint, and its impact on change detection
accuracy is examined in Section~\ref{subsec:changedetect}.

Prompts are kept \emph{identical} across all three models to ensure a controlled
comparison. An optional grounded mode first extracts readable on-screen text tokens
and then conditions the description on that evidence, reducing hallucinated detail;
in the experiments reported here the standard (ungrounded) mode was used so that
all three models operate under exactly the same conditions.

\subsection{Storage Layer}

Results are persisted to a per-run SQLite database via \texttt{invoy\_sdk/activity\_log.py}.
Each row stores the screenshot path, frame index, run label, model name, the activity
text, the change text, and two optional extracted-evidence fields. Running the same
screenshot archive with a different model simply writes a new database with a distinct
run label, allowing side-by-side comparison with no data mutation.

% ================================================================
\section{Evaluation Methodology}
\label{sec:eval}

\subsection{Dataset}

136 screenshots were captured on 16~March~2026 from a Linux workstation running a
GNOME Wayland session at 10-second intervals over approximately 23~minutes
(07:31:59--07:54:47 UTC). The session was a continuous realistic work session
involving Python development in VS~Code, a terminal, and a web browser (including
a university grade-portal visit). No held-out split is used; all 136 frames are
evaluated. Because no human activity labels are available, the 7B model is treated
throughout as a \emph{pseudo ground-truth reference}.

\subsection{Models}

Three models were evaluated. All ran on the same CUDA GPU with the
\texttt{device=``auto''} setting:

\begin{center}
\small
\begin{tabular}{@{}llrr@{}}
\toprule
\textbf{Label} & \textbf{Family} & \textbf{Params} & \textbf{VRAM} \\
\midrule
2B & Qwen2-VL    & $\sim$2B & $\sim$5\,GB \\
3B & Qwen2.5-VL  & $\sim$3B & $\sim$7\,GB \\
7B & Qwen2-VL    & $\sim$7B & $\sim$14\,GB \\
\bottomrule
\end{tabular}
\end{center}

\noindent
The 2B and 7B models belong to the same Qwen2-VL generation; the 3B model is from
the newer Qwen2.5-VL generation, offering an architecture-controlled comparison at
a lower parameter count than the 7B reference.

\subsection{Proxy Quality Metrics}

All evaluation metrics are computed purely from the text output of each model; no
images are re-examined at evaluation time.

\textbf{Token length} is the count of non-stopword tokens matching the pattern
\texttt{[A-Za-z0-9\_./:--]\{2,\}}.

\textbf{Specificity score}~$S$ measures the fraction of tokens that are concrete
application or UI identifiers: a fixed lexicon of app/tool names
(\texttt{vscode}, \texttt{chrome}, \texttt{terminal}, \texttt{github}, etc.) plus
tokens whose suffix matches a file-extension pattern
(\texttt{.py}, \texttt{.ts}, \texttt{.json}, \texttt{.md}, \ldots):
\[
  S = \frac{\text{concrete hits}}{\max(1,\; n_{\text{tokens}})} \;\in [0,1].
\]

\textbf{Hedging rate}~$H$ counts tokens expressing uncertainty
(\emph{maybe}, \emph{appears}, \emph{seems}, \emph{might}, etc.) divided by total
token count.

For change descriptions, \textbf{action-verb presence}~$V \in \{0,1\}$ flags whether
the text contains a concrete action verb (e.g.\ \emph{switch}, \emph{open},
\emph{commit}), \textbf{state-transition presence}~$T \in \{0,1\}$ flags temporal
transition language (\emph{changed}, \emph{moved from}, \emph{now}, etc.), and
\textbf{vagueness rate}~$G$ counts filler words (\emph{stuff}, \emph{things},
\emph{various}) per token.

\subsection{Proxy Scores (0--5 Scale)}

Four aggregate scores are derived from the above metrics:
\begin{align*}
  \text{Activity Fidelity}      &= 5\bigl[0.6S + 0.4(1-H)\bigr] \\
  \text{Activity Specificity}   &= 5S \\
  \text{Change Correctness}     &= 5\bigl[0.35V + 0.35T + 0.30D\bigr] \\
  \text{Change Usefulness}      &= 5\bigl[0.50D + 0.50(1-G)\bigr]
\end{align*}
where $D$ is the Change-Delta F1 score defined in Section~\ref{subsec:drift}.

\subsection{Agreement with 7B Reference}

For each frame, the output of the 2B and 3B models is compared to the 7B output using:
\textbf{TF-IDF cosine similarity} (semantic proximity under a shared vocabulary),
\textbf{ROUGE-L F1}~\cite{lin2004rouge} (longest common sub-sequence recall/precision),
and \textbf{Jaccard similarity} (token-set overlap after stopword removal).

\subsection{Change-Delta F1}
\label{subsec:drift}

The \textbf{Change-Delta F1} (CDF1) metric measures how well each model's change text
is lexically grounded in what actually changed between consecutive activity texts. For
frame $t > 0$, the \emph{activity delta} is the symmetric difference of tokenised
activity texts:
\[
  \Delta_t = \mathcal{T}(a_{t-1}) \;\triangle\; \mathcal{T}(a_t),
\]
where $\mathcal{T}(\cdot)$ is the stopword-filtered token set. Let $\mathcal{C}_t$ be
the token set of the change text at frame~$t$. Define:
\begin{align*}
  \text{precision}_t &= \tfrac{|\Delta_t \cap \mathcal{C}_t|}{|\mathcal{C}_t|} \\[4pt]
  \text{recall}_t    &= \tfrac{|\Delta_t \cap \mathcal{C}_t|}{|\Delta_t|}
\end{align*}
\[
  \text{CDF1}_t = \frac{2 \cdot \text{precision}_t \cdot \text{recall}_t}
                       {\text{precision}_t + \text{recall}_t}
\]
with $\text{CDF1}_t = 1$ if both sets are empty (correctly silent), and $0$ if exactly
one is empty. CDF1 is averaged over frames $t = 1, \ldots, 135$.

\textbf{Precision} penalises change texts that mention tokens unrelated to what
shifted. \textbf{Recall} rewards change texts that cover the actual delta tokens.
The F1 combination prevents gaming either direction alone---a verbose model cannot
inflate recall without also suffering on precision, and a terse model cannot inflate
precision without suffering on recall.

As a complementary scalar, the \textbf{Spearman rank correlation} between per-frame
activity drift $\delta_t$ and per-frame CDF1$_t$ measures whether change-text grounding
scales proportionally with the magnitude of the activity shift.

% ----------------------------------------------------------------
\subsection{Metric Provenance and Novelty}
\label{subsec:provenance}

The evaluation metrics used in this work fall into two categories: established measures
adapted from the NLP and information-retrieval literature, and task-specific proxies
introduced here.

\textbf{Established metrics.}
\emph{Token length} is a standard output-length statistic used across NLG evaluation
studies~\cite{lin2004rouge}.
\emph{Hedging rate} builds directly on prior work in computational linguistics on
hedge detection: Hyland's metadiscourse framework~\cite{hyland2005metadiscourse}
catalogues epistemic hedges (e.g.\ \textit{possibly}, \textit{seems}, \textit{might})
across academic genres, and the BioScope corpus~\cite{vincze2008bioscope} provides a
manually annotated resource for uncertainty and negation in biomedical text.
Our hedging lexicon (\textit{maybe, possibly, unclear, uncertain, likely, appears,
seems, might, could}) is adapted directly from those precedents.
\emph{TF-IDF cosine similarity}, \emph{ROUGE-L}, and \emph{Jaccard overlap} are
all standard reference-based NLG metrics; ROUGE-L in particular is
from~\cite{lin2004rouge}.

\textbf{Novel metrics.}
The \emph{specificity score} $S$ is a task-specific proxy introduced in this work.
While it is structurally inspired by precision-oriented retrieval evaluation, no prior
measure explicitly counts concrete desktop application identifiers and file-extension
tokens as a fraction of description length for activity-tracking evaluation.
We claim $S$ as an original contribution scoped to the desktop activity domain.

The \emph{Change-Delta F1} (CDF1) is fully novel.
No prior work applies symmetric-difference framing to consecutive natural-language
generation outputs to measure grounding of a change description.
The metric is structurally inspired by ROUGE-F1~\cite{lin2004rouge} and by
delta-based change detection concepts in software engineering, but the formulation ---
computing the symmetric difference of tokenised activity texts as the reference signal ---
has no direct precedent in the VLM or activity-recognition literature.

% ================================================================
\section{Results}
\label{sec:results}

\begin{table*}[t]
\centering
\caption{Aggregate evaluation results over 136 frames.
         Qwen2-VL 7B serves as pseudo ground-truth reference
         (\texttt{---} = self-comparison not applicable).
         \textbf{Bold} denotes the best value in each row.}
\label{tab:results}
\setlength{\tabcolsep}{8pt}
\begin{tabular}{@{}lrrr@{}}
\toprule
\textbf{Metric}
  & \textbf{Qwen2-VL 2B}
  & \textbf{Qwen2.5-VL 3B}
  & \textbf{Qwen2-VL 7B} \\
\midrule
\multicolumn{4}{@{}l}{\small\textit{Description quality}} \\[3pt]
Average token length
  & 37.2
  & 25.9
  & \textbf{44.7} \\
Specificity score
  & \textbf{0.038}
  & 0.016
  & 0.033 \\
Hedging rate
  & 0.013
  & 0.010
  & \textbf{0.006} \\[8pt]
\multicolumn{4}{@{}l}{\small\textit{Proxy scores (0--5 scale)}} \\[3pt]
Activity Fidelity
  & \textbf{2.09}
  & 2.03
  & 2.09 \\
Activity Specificity
  & \textbf{0.190}
  & 0.081
  & 0.165 \\
Change Correctness
  & \textbf{2.57}
  & 2.23
  & 2.15 \\
Change Usefulness
  & \textbf{3.33}
  & 3.17
  & 3.00 \\[8pt]
\multicolumn{4}{@{}l}{\small\textit{Change-text grounding}} \\[3pt]
Change-Delta F1
  & \textbf{0.334}
  & 0.269
  & 0.200 \\
Spearman $r$ (drift vs CDF1)
  & \textbf{0.863}
  & 0.814
  & 0.819 \\[8pt]
\multicolumn{4}{@{}l}{\small\textit{Agreement with 7B reference}} \\[3pt]
Cosine similarity
  & \textbf{0.338}
  & 0.327
  & ---\\
ROUGE-L F1
  & \textbf{0.292}
  & 0.285
  & --- \\
Jaccard similarity
  & \textbf{0.235}
  & 0.226
  & --- \\
\bottomrule
\end{tabular}
\end{table*}

Table~\ref{tab:results} summarises aggregate results across all 136 frames.
The following sub-sections examine each dimension in turn.

\subsection{Description Quality}
\label{subsec:desc_quality}

\begin{figure}[t]
  \centering
  \includegraphics[width=\columnwidth]{01_quality_metrics_bar.png}
  \caption{Activity description quality metrics per model (specificity and hedging
           rate on a 0--1 scale; verbosity normalised to [0,1] by dividing by the
           maximum mean token count). The 2B model achieves the highest
           specificity (0.038) despite its smaller size, while the 3B model is
           markedly more concise. Hedging rate decreases monotonically with
           parameter count.}
  \label{fig:quality}
\end{figure}

The 7B model produces the richest descriptions, averaging 44.7~tokens per
frame---20\% more than 2B (37.2) and 73\% more than 3B (25.9).
Figure~\ref{fig:quality} shows that this verbosity gap is accompanied by the lowest
hedging rate (0.006 vs.\ 0.013 for 2B), indicating that larger models express more
confident, declarative descriptions. Smaller models appear to compensate for
uncertainty by using qualified language (\emph{appears to be}, \emph{seems like}),
which in turn penalises Activity Fidelity.

A counterintuitive result is that the 2B model achieves the \emph{highest} specificity
score (0.038), surpassing 7B (0.033). This is explained by the behavioural difference
between models of different capacities: the 2B model tends to anchor descriptions on
the most visually salient keywords---application names, filenames, visible URLs---in a
surface-level extraction style, while the 7B model generates more fluent, narrative
prose that dilutes keyword density even when semantically richer. The 3B model scores
substantially below both (0.016), reflecting its tendency towards short, abstract
summaries that omit specific context.

\subsection{Proxy Score Analysis}
\label{subsec:proxy}

\begin{figure}[t]
  \centering
  \includegraphics[width=\columnwidth]{02_proxy_scores_bar.png}
  \caption{Proxy quality scores (0--5) across four evaluation dimensions.
           The 2B model leads on both change-related dimensions (Change Correctness
           and Change Usefulness) as well as Activity Specificity. Activity Fidelity
           is nearly equal for 2B and 7B. The 3B model is lowest on all four axes.}
  \label{fig:proxy}
\end{figure}

Figure~\ref{fig:proxy} shows the four proxy dimensions. Activity Fidelity is nearly
identical for 2B and 7B (both 2.09), with 3B trailing (2.03). Activity Specificity
reproduces the specificity-score ordering: 2B (0.190) $>$ 7B (0.165) $>$ 3B (0.081),
the last being less than half the 2B value.

For change descriptions, Change Correctness and Change Usefulness both incorporate
the Change-Delta F1 score~$D$ (Section~\ref{subsec:changedetect}). The 2B model
leads on both: Change Correctness 2.57 vs.\ 2.23 (3B) and 2.15 (7B); Change
Usefulness 3.33 vs.\ 3.17 (3B) and 3.00 (7B). This ordering mirrors the CDF1
ranking and is examined in Section~\ref{subsec:changedetect}. Across all four proxy
dimensions the 3B model is uniformly lowest, as confirmed by the radar chart in
Figure~\ref{fig:radar} (Section~\ref{sec:discussion}).

\subsection{Agreement with 7B Reference}
\label{subsec:agreement}

\begin{figure}[t]
  \centering
  \includegraphics[width=\columnwidth]{09_pairwise_heatmap.png}
  \caption{Mean pairwise TF-IDF cosine similarity of activity texts across
           all 136 frames. Self-similarity on the diagonal is 1.0 by definition.
           Inter-model cosines cluster in the 0.31--0.34 range, indicating
           substantial output divergence even for identical input frames.}
  \label{fig:heatmap}
\end{figure}

Figure~\ref{fig:heatmap} presents the mean pairwise cosine similarity matrix.
The 2B model is marginally closer to 7B (cosine 0.338, ROUGE-L 0.292, Jaccard 0.235)
than the 3B model (0.327, 0.285, 0.226). The 2B--3B similarity (0.315) is even lower,
indicating that the two compact models interpret the same screens in substantially
different ways. All inter-model similarities fall in the 0.31--0.34 range, which is
far below values that would indicate consistent interpretation. This reflects the
fundamentally open-ended nature of the task: with no constrained output vocabulary,
models can produce equally plausible but lexically dissimilar descriptions of the same
scene.

\subsection{Temporal and Segment Analysis}
\label{subsec:temporal}

\begin{figure}[t]
  \centering
  \includegraphics[width=\columnwidth]{s3_segment_agreement.png}
  \caption{TF-IDF cosine and Jaccard agreement of 2B and 3B with 7B, broken down
           by application context segment (keyword-based, all three models combined).
           Segments with fewer than 5 frames are omitted. The 2B model is closer to
           7B in the two longest stable browser segments (frames 63--108 and
           117--130); the 3B model is closer in shorter transitional segments.
           Segment frame counts ($n$) are shown on the horizontal axis.}
  \label{fig:timeseries}
\end{figure}

Keyword-based segmentation of the session (scanning activity texts from all three
models combined) reveals a \textbf{browser-dominant session}: approximately 80\% of
frames describe browser activity, with brief VS~Code appearances (1--2 frame bursts)
and a sustained terminal segment only in the final 6 frames (131--136). The two
known task-switch anchors (frame~$\sim$40 and frame~$\sim$110) are both recovered
as keyword transition points.

Figure~\ref{fig:timeseries} shows agreement with 7B broken down by segment. The
pattern is context-dependent: in the two longest stable browser segments (frames
63--108 and 117--130), the 2B model is consistently closer to 7B (cosine 0.400 and
0.399 vs.\ 0.325 and 0.340 for 3B). In shorter or transitional segments, the 3B
model is occasionally closer. This suggests 2B's overall agreement advantage is
driven by its stronger consistency in sustained, content-stable portions of the
session rather than being uniform throughout.

\subsection{Change Detection Selectivity}
\label{subsec:changepr}

\begin{figure}[t]
  \centering
  \includegraphics[width=\columnwidth]{s4_change_detection_pr.png}
  \caption{Precision, recall, and F1 for change detection against keyword-transition
           ground truth (segment boundaries derived from combined keyword analysis of
           all three models). All models achieve recall $\approx 1.0$ but precision
           $\approx 0.12$: they fire a change event on $\sim$130 of 136 frames, but
           only 16 frames are genuine context transitions.}
  \label{fig:changepr}
\end{figure}

Figure~\ref{fig:changepr} presents the most concrete change-detection result in
this evaluation. Using keyword-identified context boundaries (VS~Code, browser,
terminal) as ground truth, each model's change text is classified as reporting a
change if it contains an action verb or state-transition token (the V and T metrics
from Section~\ref{sec:eval}). Against 16 genuine transitions across 136 frames:

\begin{itemize}
  \item \textbf{Recall $\approx 1.0$} for all three models---every real context
        switch is captured.
  \item \textbf{Precision $\approx 0.12$}---the models fire change events on
        $\sim$130 frames, yielding $\sim$113 false positives per model.
  \item \textbf{F1 $\approx 0.21$}---far below useful task-switch detection.
\end{itemize}

The models are \emph{indiscriminate}: the V and T tokens are present in nearly
every change text regardless of whether a real transition occurred.  This result
is independent of model size---all three behave identically in terms of selectivity.
It establishes definitively that raw VLM change texts cannot serve as task-switch
signals without a downstream classifier.

\subsection{Change-Text Grounding (CDF1)}
\label{subsec:changedetect}

\begin{figure}[t]
  \centering
  \includegraphics[width=\columnwidth]{07_change_delta_f1_bar.png}
  \caption{Mean Change-Delta F1 per model (frames $t \geq 1$). The 2B model achieves
           the highest lexical grounding (0.334), meaning its change texts cover more
           of the tokens that actually appeared or disappeared between consecutive
           activity texts. The 7B model is lowest (0.200) despite producing the
           richest activity descriptions---its change texts paraphrase rather than
           echo the activity delta tokens.}
  \label{fig:alignment}
\end{figure}

Figure~\ref{fig:alignment} presents the mean CDF1 scores. The 2B model is highest
(0.334), the 3B model is intermediate (0.269), and the 7B model is lowest (0.200).
The ordering is the \emph{reverse} of the activity description richness ranking,
which makes it the most interpretively interesting result in the evaluation.

The explanation lies in how each model generates change text. The 2B model's
surface-level, keyword-anchored activity descriptions produce activity deltas that
are small and dominated by specific tokens (app names, filenames). Its change texts
tend to copy the same identifiers directly, yielding high overlap. The 7B model
generates fluent, narrative activity texts; when these shift between frames, the
delta is large and spread across many paraphrased concepts. The 7B change texts
describe the shift in natural language that is semantically accurate but uses
different vocabulary---which the CDF1 metric penalises because CDF1 is lexical, not
semantic. This is a known trade-off in token-overlap-based evaluation.

Figure~\ref{fig:scatter} reveals the mechanism behind the CDF1 gap via OLS
regression. All three intercepts are near-identical ($\leq$0.017)---models do not
differ in their baseline grounding at near-zero drift. What differs is the
\emph{slope}: 2B=0.697, 3B=0.545, 7B=0.394. The slope gap between 2B and 7B
(0.303, bootstrap 95\% CI [0.256, 0.349]) excludes zero, confirming that 2B's
change-text grounding rises significantly faster with activity drift. In practical
terms: 2B benefits most from high-drift frames (large activity text shifts), where
its keyword-anchored style produces tightly matched change texts, while 7B's
paraphrasing style yields smaller gains at the same drift magnitude.

\begin{figure*}[t]
  \centering
  \includegraphics[width=\textwidth]{07b_drift_vs_cdf1_scatter.png}
  \caption{Per-frame scatter of activity drift vs Change-Delta F1 for each model,
           with OLS regression line. Spearman rank correlations: 0.863 (2B), 0.814
           (3B), 0.819 (7B). OLS slopes differ significantly: 2B=0.697, 3B=0.545,
           7B=0.394 (slope CI for 2B$-$7B: [0.256, 0.349], excludes zero). Intercepts
           are near-identical ($\leq$0.017). The difference is therefore in
           \emph{slope}: 2B's grounding rises faster with drift than 7B's, not from
           a higher baseline.}
  \label{fig:scatter}
\end{figure*}

% ================================================================
\section{Discussion}
\label{sec:discussion}

\begin{figure}[t]
  \centering
  \includegraphics[width=\columnwidth]{08_radar_overall.png}
  \caption{Radar chart across six normalised evaluation dimensions. The 2B model
           (blue) leads on the change-related axes (Change-Delta F1, Change
           Correctness, Change Usefulness) while the 7B model (green) leads on
           activity description richness (token length, Activity Fidelity). The 3B
           model (orange) is interior to both competitors on every axis.}
  \label{fig:radar}
\end{figure}

\subsection{Model Size vs.\ Architecture (RQ1, RQ2)}

The results support a \textbf{parameter-count-first} interpretation for
\emph{activity description} quality (RQ1): the 7B model leads on richness (44.7 tokens),
confidence (hedging rate 0.006), and Activity Fidelity. For \emph{change-text} quality
the ranking reverses---the 2B model leads on CDF1, Change Correctness, and Change
Usefulness---a distinction examined in Section~\ref{subsec:changedetect}.

RQ2 is answered negatively: the newer Qwen2.5-VL architecture does not compensate
for reduced parameter count at the 3B scale. Across all eleven metrics in
Table~\ref{tab:results}, the 3B model scores lower than either the 2B or 7B
Qwen2-VL model. The 3B model's descriptions are systematically shorter, more
abstract, and lower in concrete identifier density. Figure~\ref{fig:radar} makes
this clear in a single view: the 3B polygon is interior to both competitors on
every axis. A plausible explanation is that the architectural improvements in
Qwen2.5-VL---enhanced spatial understanding, improved OCR alignment---primarily
manifest at 7B-scale and above, where they can be fully utilised; at 3B the
capacity constraint becomes the binding factor.

The 2B and 7B models from the same Qwen2-VL generation show a marginally higher
mean cosine (0.338 vs.\ 0.327 for 3B--7B). However, a Wilcoxon signed-rank test
on per-frame cosine similarities yields $p = 0.104$---not statistically significant
at $\alpha = 0.05$ (2B is closer to 7B on only 54.4\% of frames). The mean
difference is noted for completeness but is not claimed as a robust architectural
finding.

\subsection{Change-Text Grounding and RQ3}

The CDF1 result---2B leads, 7B is last---requires careful interpretation. CDF1 is a
\emph{lexical} metric: it rewards change texts that reuse the same tokens present in
the activity delta. A model that writes ``VS~Code was replaced by Chrome'' scores
higher than one that writes ``the user transitioned from a development environment to
a web browser'', even though both convey the same meaning.

This means 7B's low CDF1 is not evidence that it produces poor change texts.
Rather, it reflects that 7B generates more fluent, semantically rich language that
paraphrases the delta instead of echoing it verbatim. For downstream consumers that
parse the change text by keyword matching (e.g.\ extracting which app was opened),
2B's lexical grounding is directly useful. For consumers that use embedding similarity
or LLM-based comparison, 7B's richer paraphrases may be equally or more useful.

Crucially, all three models show strong Spearman correlation between activity drift
and CDF1 ($r = 0.863$, $0.814$, $0.819$; Figure~\ref{fig:scatter}). This means every
model \emph{scales} its change-text effort proportionally with the magnitude of the
activity shift---the difference is only in the baseline vocabulary choice.

RQ3 is answered in two parts. For \emph{selectivity}, the precision/recall
evaluation (Section~\ref{subsec:changepr}) is unambiguous: all three models have
precision $\approx 0.12$ and recall $\approx 1.0$, producing $\sim$113 false
positives per model against 16 real transitions. Change texts are not selective---
they describe differences on almost every frame regardless of whether a real context
switch occurred. For \emph{grounding quality} on genuine change frames, the CDF1
metric differentiates the models (2B leads), but this only matters after a
post-processing layer that can distinguish true transitions from the continuous stream
of over-triggered change texts. Whether 2B's lexical grounding or 7B's semantic
paraphrasing is preferable for that downstream step depends on the consumer's
matching strategy.
We identify three complementary directions for future work:
\begin{enumerate}
  \item \textbf{Temporal classifier.} Train a lightweight binary classifier on
        pairs of consecutive activity texts to distinguish continuations from
        task switches, using the richer 7B descriptions as pseudo-labels.
  \item \textbf{Semantic CDF1.} Replace token-overlap with embedding cosine
        similarity between the change text and the activity delta, so that
        semantically equivalent paraphrases (7B-style) are credited alongside
        verbatim echoes (2B-style).
  \item \textbf{Fine-tuning on labelled transitions.} Even a small set of
        human-annotated task-switch frames---perhaps 200--300 pairs---could provide
        sufficient supervision for a specialised change-detection head.
\end{enumerate}

\subsection{Limitations}

Four limitations bound the scope of these conclusions:
\begin{enumerate}
  \item \textbf{No human ground truth.} All metrics are proxy-based. Specificity
        heuristics reward keyword density, not semantic accuracy; a model that
        invents plausible-looking filenames would score well without being correct.
  \item \textbf{Single session.} One user, one machine, one 23-minute window.
        Task diversity is limited to coding, terminal work, and browsing; models
        may rank differently on email composition, document editing, or video
        playback.
  \item \textbf{CDF1 is lexical, not semantic.} Change-Delta F1 rewards token-level
        overlap between the change text and the activity delta. A semantically
        accurate paraphrase scores lower than a verbatim echo. Models with richer,
        more fluent language (7B) are structurally penalised relative to models with
        surface-level extraction styles (2B), even when semantic content is equivalent.
  \item \textbf{No latency or throughput measurement.} Production deployment
        requires characterising inference time per frame; at these model sizes the
        effective capture cadence is bounded by inference, not by the configured
        interval.
\end{enumerate}

\subsection{Deployment Guidance}

Based on the results, we offer the following practical recommendations:
\begin{itemize}
  \item \textbf{Use 7B} when GPU resources allow ($\geq$14\,GB VRAM), description
        fluency matters, and the change-text consumer uses semantic comparison
        (embeddings, LLM evaluation). The 7B model produces the richest activity
        descriptions despite lower lexical grounding in change texts.
  \item \textbf{Use 2B} when memory is constrained ($\sim$5\,GB VRAM) \emph{or}
        when the change-text consumer uses keyword matching or structured parsing.
        The 2B model leads on Change-Delta F1 (0.334) and Spearman calibration
        (0.863), making its change texts the most lexically grounded.
  \item \textbf{Avoid 3B} unless output brevity is the explicit design goal. For
        every quality dimension measured, 2B is a superior choice at comparable
        or lower VRAM cost.
  \item \textbf{Post-process change texts} before using them as task-switch signals.
        Even the highest-CDF1 model (2B at 0.334) leaves most delta tokens uncovered;
        a structured classifier downstream is necessary for reliable task-switch
        detection.
\end{itemize}

% ================================================================
\section{Related Work}
\label{sec:related}

\paragraph{Desktop activity recognition.}
Early passive monitoring systems used application-level event logs or window-manager
APIs to infer activity, covering only the applications they were explicitly
instrumented for. Vision-based approaches---including screen-region classifiers and
OCR-assisted labelling---offered broader coverage but were fragile to unseen
applications and required domain-specific training data. Our approach replaces
hand-crafted classifiers with a generalist VLM, enabling open-vocabulary activity
descriptions with no application-specific engineering.

\paragraph{Vision-language models for UI understanding.}
Qwen2-VL~\cite{wang2024qwen2vl} and Qwen2.5-VL~\cite{qwen25vl2025} represent the
state of the art in open-weight vision-language models optimised for document and
interface understanding, with explicit OCR alignment and spatial grounding in the
training objective. Related work on MiniGPT-4~\cite{zhu2023minigpt4} and similar
models established the feasibility of instruction-following on visual inputs at
sub-10B scales, motivating the localdeployment approach taken here.

\paragraph{Proxy evaluation for generative models.}
ROUGE-L~\cite{lin2004rouge} and TF-IDF cosine similarity are standard automatic
metrics for text quality when reference texts are available but human annotation
is prohibitively expensive. Their known limitations---sensitivity to phrasing over
semantics---motivate the additional specificity and hedging heuristics we introduce,
which reward grounding in visible UI elements rather than fluent but
content-neutral prose.

% ================================================================
\section{Conclusion}
\label{sec:conclusion}

We presented Invoy, an SDK for continuous desktop activity tracking using local
vision-language models, and reported the first systematic comparison of three
open-weight VLMs---Qwen2-VL 2B, Qwen2.5-VL 3B, and Qwen2-VL 7B---on a real-world
screenshot sequence. Our findings are threefold. First, parameter count governs \emph{activity description}
quality: the 7B model leads on richness, fluency, and confidence. Second, the newer
Qwen2.5-VL architecture does not compensate for the parameter deficit at 3B, making
the older Qwen2-VL 2B the superior compact choice across all dimensions. Third, for
\emph{change-text grounding}, the ranking reverses: the 2B model achieves the highest
Change-Delta F1 (0.334 vs.\ 0.200 for 7B), because its keyword-anchored activity
texts produce lexically tight deltas that its change texts echo directly. This
finding highlights a practical split: use 7B when the downstream consumer interprets
change texts semantically, and 2B when keyword-level grounding is required.
A structured classifier downstream is essential: all three models achieve recall
$\approx 1.0$ but precision $\approx 0.12$ against keyword-identified transitions,
meaning $\sim$130 of 136 frames trigger a change text regardless of actual context
continuity.

The Invoy SDK, evaluation scripts, and generated databases are available in the
project repository.

% ================================================================
\begin{thebibliography}{6}

\bibitem{wang2024qwen2vl}
J.~Wang et al.,
``Qwen2-VL: Enhancing Vision-Language Model's Perception of the World at Any
  Resolution,''
\textit{arXiv preprint arXiv:2409.12191}, 2024.

\bibitem{qwen25vl2025}
Qwen Team, Alibaba Group,
``Qwen2.5-VL Technical Report,''
\textit{arXiv preprint arXiv:2502.13923}, 2025.


\bibitem{zhu2023minigpt4}
D.~Zhu et al.,
``MiniGPT-4: Enhancing Vision-Language Understanding with Advanced Large
  Language Models,''
\textit{arXiv preprint arXiv:2304.10592}, 2023.

\bibitem{lin2004rouge}
C.-Y.~Lin,
``ROUGE: A Package for Automatic Evaluation of Summaries,''
in \textit{Proc.\ ACL Workshop on Text Summarization Branches Out}, 2004.

\bibitem{hyland2005metadiscourse}
K.~Hyland,
\textit{Metadiscourse: Exploring Interaction in Writing}.
Continuum, 2005.

\bibitem{vincze2008bioscope}
V.~Vincze, G.~Szarvas, R.~Farkas, G.~M\'{o}ra, and J.~Csirik,
``The BioScope Corpus: Biomedical Texts Annotated for Uncertainty,
  Negation and Their Scopes,''
in \textit{Proc.\ BioNLP Workshop at ACL}, 2008, pp.~38--45.

\end{thebibliography}

\end{document}
